\paragraph{Định nghĩa và cơ sở toán học:}
\begin{itemize}
    \item Hệ số tương quan ($\rho$) giữa hai biến ngẫu nhiên $X_1$ và $X_2$ được tính bằng tỷ số giữa hiệp phương sai và tích các độ lệch chuẩn của chúng:
    \bea
    \rho(X_1, X_2) = \frac{\text{cov}(X_1, X_2)}{\sqrt{\text{var}(X_1)\text{var}(X_2)}} .
    \eea
    \item Đối với một vectơ ngẫu nhiên đa biến $X$, ma trận tương quan $P = \rho(X)$ có các phần tử tại vị trí $(i,j)$ là tương quan cặp giữa $X_i$ và $X_j$. Ma trận này luôn có tính chất xác định dương hoặc ít nhất là nửa xác định dương.
\end{itemize}

\paragraph{Các đặc tính chính:}
\begin{itemize}
    \item Hệ số tương quan tuyến tính luôn nằm trong khoảng $[-1, 1]$. Khi $\rho = 1$, hai biến có sự phụ thuộc tuyến tính thuận hoàn hảo, và khi $\rho = -1$, hai biến có sự phụ thuộc tuyến tính nghịch hoàn hảo.
    \item Một thuộc tính quan trọng khác là tính bất biến: tương quan tuyến tính có tính bất biến dưới các phép biến đổi tuyến tính tăng nghiêm ngặt. Tuy nhiên, nó không bất biến đối với các phép biến đổi phi tuyến.
\end{itemize}

\paragraph{Hạn chế và những cạm bẫy:}
Tương quan chỉ được xác định khi các biến có phương sai hữu hạn, điều này gây khó khăn với các rủi ro có đuôi nặng (như rủi ro hoạt động) có phương sai vô hạn. Mối liên hệ giữa tương quan và độc lập cũng dễ gây nhầm lẫn: độc lập dẫn đến tương quan bằng $0$, nhưng ngược lại không luôn đúng. Các biến không tương quan vẫn có thể phụ thuộc chặt chẽ về mặt cấu trúc trong các mô hình phi Gaussian.\\
Bên cạnh đó, mô hình này còn vấp phải ba cạm bẫy lớn trong quản trị rủi ro:
\begin{itemize}
    \item \textbf{Cạm bẫy 1 (Xác định phân phối):} Các phân phối biên và tương quan cặp không đủ để xác định phân phối đồng thời của một vectơ ngẫu nhiên.
    \item \textbf{Cạm bẫy 2 (Tương quan có thể đạt được):} Với các phân phối biên cho trước, dải tương quan đạt được thực tế có thể hẹp hơn khoảng $[-1, 1]$. Ngay cả khi hai biến phụ thuộc thuận mạnh nhất, tương quan tuyến tính vẫn có thể rất thấp.
    \item \textbf{Cạm bẫy 3 (Rủi ro danh mục):} Một quan niệm sai lầm phổ biến là cho rằng rủi ro của một danh mục là tệ nhất khi tương quan đạt cực đại. Thực tế, đối với các phân phối không đối xứng, danh mục có tương quan thấp hơn đôi khi lại tạo ra tổn thất cực đoan lớn hơn so với tương quan hoàn hảo.
\end{itemize}

\noindent Mặc dù tương quan tuyến tính là một tham số tự nhiên trong các mô hình elip, nó không phản ánh được các cấu trúc phụ thuộc phi tuyến hoặc phụ thuộc đuôi. Do đó, các nhà quản lý rủi ro thường hướng tới sử dụng Copula để tách biệt hành vi biến của tài sản khỏi cấu trúc phụ thuộc của chúng, cho phép đánh giá chính xác hơn các kịch bản rủi ro cực đoan.
